% Standalone addendum fragment. No new package dependencies.
\subsection{The uniform degree-ten Paley net has no rational integral member}
\label{app:paley-net-rational-exclusion}

We work over an algebraically closed field of characteristic zero. Put
$\eta^2+\eta+2=0$, $\alpha=(\eta-1)/2$, and $\gamma=3(\eta+3)/4$.
On $\mathbb F_3$, with negative section $C_0$ and fiber $f$, consider the
net $V\subset H^0(10C_0+34f)$ whose members have multiplicity at least
four at $Q_i=(\zeta^i,\zeta^{5i})$ and at least six at
$T_i=(\alpha\zeta^i,\gamma\zeta^{5i})$, $0\le i<7$. Here $\zeta$ is a primitive
seventh root. Its dimension is three. An exact basis $F_0,F_1,F_2$ has
characters $1,3,5$ under $(X,Y)\mapsto(\zeta X,\zeta^5Y)$, normalized
by the coefficients $25,27,20$ of $Y^{10},X^2Y^{10},X^4Y^{10}$,
respectively. These specifications and the prescribed jets determine the
basis over $\mathbb Q(\eta)$.

\begin{theorem}
\label{thm:paley-net-no-rational}
No member of this net is an integral curve with rational normalization.
In particular, the Paley uniform degree-ten, four-pole norm net cannot
supply an augmentation through an integral norm equation. This statement concerns this net, not all
possible eight-word constructions or all degree-ten incidence patterns.
\end{theorem}

\begin{proof}
We give the geometric argument and specify the finite certificates on
which it depends. Reduction is at the prime above $29$ with $\eta=7$;
thus $\alpha=3$ and $\gamma=22$. The exact basis and its reduction are checked
coefficientwise, using two orbit representatives for the prescribed jets.
The $217$ prescribed jet conditions on $220$ coefficients have rank
$217$ modulo $29$. For each of the three indicated characters, the
$31$ representative jet equations in its $32$ allowed monomials give
one exact eigenvector over $\mathbb Q(\eta)$. The independently verified
vectors have the stated normalizations and reduce to the modular kernel.
The unit rank minor therefore proves that they span the entire net and
its integral model. Every nonzero member has actual class
$10C_0+34f$: the leading coefficient is $25a+27bX^2+20cX^4$, and the
three weighted-infinity restrictions below have disjoint supports.

Besides the fourteen prescribed basepoints, the net has exactly two
simple basepoints: $(X,Y)=(0,0)$ and $(t,Z)=(0,0)$ in the chart
$t=1/X$, $Z=Y/X^3$. One blowup at each of these sixteen points resolves
the net. On the resulting smooth surface $S$, its moving line bundle $L$
satisfies
\[
 L^2=14,\qquad K_S L=14,\qquad c_2(S)=20,\qquad p_a(L)=15.
\]
These assertions hold over the DVR at the chosen prime: the closed
relative base locus has empty special fiber, and properness then excludes
basepoints on the geometric generic fiber.

Here are the boundary checks needed for the discriminant computation.
On the affine chart, the nine two-by-two minors of the jet matrix have
no common divisorial factor. Along the fourteen exceptional curves, the
net restrictions are basepoint-free tangent-cone maps of degree at most
six. Their coefficient matrices have rank three; the displayed
three-column minors at the two orbit representatives are $5$ and $20$
modulo $29$. The two remaining exceptional curves have degree-one maps.
On $C_0$ the restriction is $[25:27X^2:20X^4]$, and on the original
infinity fiber it is represented by
\[
 [23Z^8+10Z:\;15Z^9+7Z^2:\;20Z^{10}+2Z^3].
\]
After removal of the common basepoint factor this is nonconstant of
degree at most ten. All these degrees are below $29$, so the maps are
generically separable. Thus the jet matrix has rank at least two at the
generic point of every boundary divisor. Its rank-one locus is finite
on all of $S$.

The universal first-jet incidence on $S\times\mathbb P^2$ consequently
has the expected dimension one: over jet-rank-two points its fiber is a
point, and over the finitely many rank-one points its fiber is a line.
Jet rank zero is excluded by basepoint-freeness. Its cycle is therefore
the top Chern class, and its pushforward is an effective discriminant
divisor of degree
\[
 c_2(J^1L)=c_2(S)+2K_SL+3L^2=90.
\]
This is a statement about a divisor with multiplicities; it assumes
neither generic nodality nor that every component has Gauss degree one.

The seven old cubic graphs yield seven parameter lines. Each contributes
multiplicity four. Indeed, the ramification determinant is simple along
each old graph; the two normal sections on the first graph are the
coprime quartics
\[
 24+14X+2X^2+27X^3+12X^4,\qquad
 25+27X+28X^2+9X^3+6X^4.
\]
They define a separable degree-four map to its parameter line.
Equivariance supplies the other six cases. Subtracting these contributions
leaves an effective discriminant divisor of degree $62$.

Modulo $29$, the remaining ramification factor $\Gamma$ has $135$ terms
and bidegree $(73,21)$. It is irreducible, and its smooth rational point
$(1,10)$, with gradient $(27,14)$, proves geometric irreducibility.
There are $4,018$ distinct certified Gauss-image points over $\mathbb
F_{29^3}$. Each receipt includes a point of $\Gamma$, a rank-two jet
witness, and its normalized parameter image. Interpolation in all seven
character spaces of homogeneous degree $62$ has total kernel dimension
one: the character-five block has rank $287$ of $288$, and the other six
blocks have full rank. Its equation $H(a,b,c)$ has $271$ terms.
Since $4,018>62^2$, B\'ezout's theorem implies that every interpolated
relation contains the Gauss image. If that image had degree less than
$62$, all its degree-$62$ multiples would give a kernel of dimension
larger than one. Thus the image has degree $62$ and exhausts the residual
discriminant cycle, including possible contributions from isolated
rank-one points.

For completeness, this identifies the reduction of a characteristic-zero
divisor, not just a modular curve. On the smooth relative surface times
the parameter plane, the three jet equations and the uniformizer form a
regular sequence, since the special incidence has expected codimension.
The incidence is therefore flat over the DVR, and proper pushforward of
its fundamental cycle commutes with specialization. Subtract the seven
horizontal old-line contributions. The primitive equation of the
remaining degree-$62$ divisor reduces to $H$ up to a nonzero scalar.

An integral rational original member must have multiplicity at least
$15$ in this residual divisor. One minor boundary issue is relevant here.
The rank-three tangent matrices prohibit multiplicity jumps at the
fourteen prescribed points. At the two additional simple points the
linear-part kernel is exactly $[0:1:0]$. The corresponding member $F_1$
has order exactly two at each, with quadratic terms $24Y^2$ and $7Z^2$
modulo $29$. Thus its resolved divisor is reduced as well: if the
original member were integral rational, its components would be its
rational strict transform and the two exceptional curves, each once.
For any of these reduced connected divisors with $k$ rational components,
\[
 \sum\delta=p_a(L)+k-1\ge15.
\]
A generic pencil through it avoids its singular points with its other
section, so its local total spaces are smooth. The discriminant
intersection length is the sum of the local Milnor numbers, at least
$\sum\delta$. Hence its parameter has discriminant multiplicity at
least $15$. An integral original member lies on none of the seven
old-factor lines, so the same assertion holds for the residual divisor.

Finally explicit polynomial identities exclude every multiplicity-$15$
point of $H$ over $\overline{\mathbb F}_{29}$. On $a=1$, a recorded linear
combination of monomial multiples of Hasse derivatives of orders below
$15$ is $1$. On $a=0,c=1$, a univariate B\'ezout identity in $b$ gives
$1$ using restrictions of those derivatives. The remaining point
$[0:1:0]$ has multiplicity six. These are coefficient identities, not
pointwise finite-field tests, and have been independently replayed.
The multiplicity-$15$ locus is closed in the proper relative parameter
plane. Its empty geometric special fiber therefore implies an empty
geometric generic fiber, even when a putative characteristic-zero
parameter specializes onto an old-factor line. This contradicts the
necessary multiplicity established above.
\end{proof}

\paragraph{Certificate provenance.}
The reproducibility directory is
\path{research/prime_line_strengthening/degree_ten_four_pole_route/}.
The exact kernel and independent coefficient check are
\path{exact_kernel.json} and \path{verify_exact_kernel.py}.
In its \path{surface/} subdirectory, the base-locus, boundary, and Gauss
receipts are \path{base_locus.json},
\path{base_locus_boundary.verified.json}, and
\path{gauss_determinant.json}; the boundary supplement is regenerated
by \path{verify_base_boundary.py}.
The point and interpolation receipts are
\path{gauss_samples.json} and
\path{gauss_image_interpolation.json}, with independent source replay
\path{verify_gauss_samples.py}.
The final explicit identities are
\path{hasse_multiplicity_gate.json}, replayed by
\path{verify_hasse_gate.py} without elimination.
The relative-cycle and exceptional-component audit is recorded in
\path{RELATIVE_DISCRIMINANT_AUDIT.md}.

\paragraph{Nonprimitive norm powers.}
The theorem alone does not exclude a norm equation that is a power
of an irreducible rational curve. This possibility is excluded by separate
small linear certificates. If a nonzero polynomial $G^t$ belongs to the
uniform norm net, where $t\mid10$ and $t>1$, then
\[
 \deg_YG\le10/t,\qquad \operatorname{wdeg}_{(1,3)}G\le\lfloor34/t\rfloor,
 \qquad \operatorname{mult}_{Q_i}G\ge\lceil4/t\rceil,
 \quad \operatorname{mult}_{T_i}G\ge\lceil6/t\rceil.
\]
The corresponding jet matrices modulo $29$ have full column rank:
\[
\begin{array}{c|c|c|c}
 t&\text{matrix size}&\text{rank}&\text{certified maximal minor}\\\hline
 2&63\mathbin{\times}63&63&27\\
 5&28\mathbin{\times}12&12&27\\
 10&14\mathbin{\times}5&5&8
\end{array}
\]
All entries are Hasse evaluations at the same fourteen points. These
nonzero minors exclude the corresponding characteristic-zero kernels.
Thus the net contains no such power, without assuming that a declared
four-pole allowance is attained. For an augmentation whose eliminated
norm equation is a scalar multiple of a pure power of its irreducible
image equation, these three cases account for every nonbirational
intermediate degree. The receipts, including selected original rows,
are \path{surface/nonprimitive_norm_gate.py} and
\path{surface/nonprimitive_norm_gate.json}.
For a proper augmentation, separability of the selected fibers and
regularity of the added function there ensure that taking the primitive
part preserves these contacts. Norm transitivity gives a power of the
minimal monic equation; Gauss's lemma gives the same power after primitive
polynomial clearing. The weighted norm bound must also hold, as in the
stated uniform four-pole construction.

\paragraph{The orbit-2 net.}
The same method closes the other rigid seven-cubic bank, but its base
geometry and interpolation have to be checked separately. Let $V_2$ be
the space of weighted-degree-at-most-$34$, $Y$-degree-at-most-$10$
polynomials having multiplicity at least four at the seven
four-agreement points and at least six at the seven three-agreement
points of the orbit-2 bank. Use its archived affine model over
$\mathbb Q(q)$, where $q^3-10q^2+3q+1=0$.

\begin{theorem}
\label{thm:orbit-two-net-no-rational}
No member of $V_2$ is an integral curve with rational normalization.
Moreover, $V_2$ contains no nonzero polynomial $G^t$ with
$t\in\{2,5,10\}$. Consequently this bank also admits no proper uniform
degree-ten four-pole augmentation whose primitive norm equation has the
stated caps and contacts.
\end{theorem}

\begin{proof}
Reduce at $q=4$ above $83$. The fourteen points remain distinct. The
$217$ prescribed Hasse equations have a $217$-column minor equal to
$3$ modulo $83$ among the $220$ columns, independently checked by an
exact integer determinant. Since there are exactly $217$ equations, this proves
rank $217$ over the exact cubic field as well. Over the chosen DVR,
solving those pivot variables with each of the three free variables
set to one gives an integral kernel basis reducing to the saved modular
basis. Thus no unverified lifting of a modular kernel is needed.

In this case there are exactly fourteen basepoints, not sixteen.
After dividing two pairwise resultants by their forced degree-$364$
factor, their residual gcd is one. Specialized ordinate gcds at every
known base coordinate contain only the prescribed point and power.
The leading and infinity restrictions have gcd one, including the
projective corner. All fourteen tangent-cone triples have rank three
and no common projective zero. Hence one blowup at each point resolves
the net, giving
\[
 L_2^2=16,\qquad K_SL_2=12,\qquad c_2(S)=18,
 \qquad p_a(L_2)=15.
\]
The rank-three tangent maps also prohibit every base-multiplicity jump:
an integral original member remains integral on the resolved surface.
The affine jet-minor gcd is one. The boundary restrictions are
nonconstant of degrees at most $4$ and $10$, and the exceptional
restrictions have degrees $4$ or $6$. All are below $83$, so they are
generically separable. The rank-one jet locus is finite on the whole
surface.

The expected-dimensional jet incidence therefore pushes to a degree-$90$
discriminant cycle. Every old graph again contributes four: the
ramification determinant is simple along it, and its two normal quartics
are coprime of maximum degree four. The residual effective cycle has
degree $62$. The remaining ramification factor $\Gamma_2$ is irreducible
modulo $83$, of bidegree $(76,21)$ with $986$ terms. Its smooth rational
point $(1,75)$, with gradient $(64,11)$, proves geometric irreducibility.

No cyclic symmetry is assumed in the interpolation. There are $4,102$
distinct certified images over
$\mathbb F_{83}[r]/(r^3-r-3)$. Full homogeneous-degree-$62$
interpolation has rank $2,015$ on $2,016$ columns; its unique relation
$H_2$ has $1,990$ terms and vanishes at all the certified points.
Since $4,102>62^2$, the same B\'ezout and kernel-dimension argument shows
that its image has degree $62$ and exhausts the residual cycle with
multiplicity one. The relative complete-intersection and proper
cycle-specialization argument in the preceding proof identifies $H_2$
as the reduction of the characteristic-zero residual discriminant.

An integral rational member would have total delta $15$, hence residual
discriminant multiplicity at least $15$, by the same generic-pencil
Milnor-number argument. It lies off the old-factor lines. Explicit
Hasse identities exclude that multiplicity everywhere in the parameter
plane: on $a=1$ a linear combination of degree-bounded derivative
multiples is $1$; on $a=0,c=1$ a univariate B\'ezout identity is $1$;
and $[0:1:0]$ is not on $H_2$. The affine certificate comes from a
full-rank $2,560$-by-$2,016$ row span and all identities are independently
replayed by coefficient summation. Properness of the closed relative
multiplicity locus then excludes the characteristic-zero member.

For the nonprimitive cases, the same smaller jet matrices as above have
full column rank modulo $83$:
\[
\begin{array}{c|c|c|c}
 t&\text{matrix size}&\text{rank}&\text{certified maximal minor}\\\hline
 2&63\mathbin{\times}63&63&67\\
 5&28\mathbin{\times}12&12&49\\
 10&14\mathbin{\times}5&5&31
\end{array}
\]
The selected original-row determinants were independently recomputed
by integer Bareiss elimination. These unit minors exclude all three
exact kernels. The primitive norm-power argument and properness
hypotheses from the preceding paragraph now give the augmentation
conclusion, without assuming that its pole allowance is attained.
\end{proof}

The separate receipts are in
\path{research/prime_line_strengthening/degree_ten_four_pole_route/orbit2_surface/}.
The complete proof and scope are in \path{EXCLUSION.md};
\path{profile.json} and \path{old_graph_degrees.json} certify the base
geometry and the seven degree-four contributions.
\path{gauss_samples.verified.json} replays every source and jet-kernel
witness; \path{gauss_image_interpolation.json} records the full
interpolation. The explicit final identities and their independent
replay are \path{hasse_multiplicity_gate.json} and
\path{hasse_multiplicity_gate.verified.json}.
The norm-power determinants are in \path{nonprimitive_norm_gate.json}
and \path{nonprimitive_norm_gate.verified.json}.
